\documentclass[]{report}
\usepackage[utf8]{inputenc} % Input UTF8
\usepackage[portuguese]{babel} %Usar língua portuguesa
\usepackage[T1]{fontenc} % Fontes T1

\usepackage{indentfirst}

\title{\textbf{Introdução à Engenharia de Computadores e Informática \\ Aula 1 de \LaTeX}}
\author{Santiago Ferreira de Almeida 139931}
\date{Universidade de Aveiro \today}

\begin{document}
\maketitle
\renewcommand{\contentsname}{Índice}
\tableofcontents

chapter{caracteres especias}\\
section{1} \textbackslash \$ \~\\
section{2} \{ \} \\
section{3} \[ \]\\
section{tipos de texto}\\
chapter{translineação}\\

Hello World from \LaTeX!\\
Ola sou o santiago sou do primeiro ano de engenharia de computadores e informatica.\\
chapter{dimensão das letras}\\
Posso {\small mudar} de {\huge tamanho} de letra sempre que
me apetecer. Essa \Large mudança permanece em efeito
{\normalsize até} ao próximo comando de mudança de tamanho,
caso exista,\normalsize\space ou até ao fim do bloco
(delimitado por \{ e \}) onde está inserida\\

chapter{lista de itens}\\
section{listas não numeradas}\\
\begin{itemize}
\item este é um item;
\item este é outro item;
\item etc.
\end{itemize}\\
\begin{itemize}
\item[+] este é um item;
\item[+] este é outro item;
\item[+] etc.
\end{itemize}\\
\begin{description}
\item[UA] - Univ. de Aveiro;
\item[UC] - Univ. de Coimbra;
\item[UP] - Univ. do Porto.
\end{description}\\
section{listas numeradas}\\
\begin{enumerate}
\item este é um item;
\item este é outro item;
\item etc.
\end{enumerate}\\
\renewcommand{\theenumi}{\Roman{enumi}}
\begin{enumerate}
\item Este é um item;
\item Este é outro item;
\item etc.
\end{enumerate}\\
\renewcommand{\theenumi}{\Roman{enumi}}
\begin{enumerate}
\item Este é um item, com subitems:
\renewcommand{\theenumi}{\alpha{enumi}}
\begin{enumerate}
\item Este é um item;
\item Este é outro item;
\item etc.
\end{enumerate}
\renewcommand{\theenumi}{\Roman{enumi}}
\item Este é outro item;
\item etc.
\end{enumerate}\\

section{figuras}\\
\begin{figure}[h]
\begin{center}
\fbox{Conteúdo da figura, pode ser texto, imagem, etc.}
\end{center}
\caption{Legenda da figura}
\end{figure}\\
\begin{figure}[t]
\begin{center}
\fbox{Conteúdo da figura, pode ser texto, imagem, etc.}
\end{center}
\caption{Legenda da figura}
\end{figure}\\
\begin{figure}[b]
\begin{center}
\fbox{Conteúdo da figura, pode ser texto, imagem, etc.}
\end{center}
\caption{Legenda da figura}
\end{figure}\\
\begin{figure}[p]
\begin{center}
\fbox{Conteúdo da figura, pode ser texto, imagem, etc.}
\end{center}
\caption{Legenda da figura}
\end{figure}\\

section{tabelas}\\
section{imagens}\\
section{matemática}\\




\textbackslashvec o.

\\ 
\sqrt[3]{x}
\text % .


\end{document}
